\bilingualchapterstar{Glossary}{术语表}

\begin{englishblock} This section preserves the original English glossary structure so that terminology, page references and lookups stay aligned with the rest of the book. The title has been made bilingual to match the bilingual edition. \end{englishblock}

\begin{chineseblock} \noindent 本节保留原书英文术语与定义结构，便于和正文中的英文术语、页码与引用保持一致；标题已切换为双语入口，可作为正文双语内容的术语检索补充。\end{chineseblock}

\medskip

\begin{chinesetableblock}
{\small
\setlength{\LTpre}{0pt} \setlength{\LTpost}{0pt} \setlength{\tabcolsep}{4pt} \renewcommand{\arraystretch}{1.1} \begin{longtable}{@{}>{\raggedright\arraybackslash\bfseries} p{0.27\textwidth}>{\raggedright\arraybackslash} p{0.69\textwidth}@{}} Active fund, Active management & A collective fund where securities are bought and sold actively to try and outperform the general market. Normally more expensive than passive funds, as active fund managers think they can generate, and charge for, alpha. Hedge funds are an extreme example of active funds. As distinct from passive funds and passive management. See page 106. \par 主动基金、主动管理：一种通过主动买卖证券来力图跑赢整体市场的集合投资基金。由于主动基金经理认为自己能够创造 alpha 并据此收费，这类基金通常比被动基金更贵。对冲基金是主动基金的极端例子。与被动基金和被动管理相对。见第 106 页。 \\[0.65em]
Alpha & The returns of an active manager or trader can be split into beta and alpha. The beta are the returns you could get from investing in the general market, i.e. in a passive index fund. Any additional return due to the manager’s skill is alpha. \par Alpha：主动管理者或交易者的收益可以拆分为 beta 和 alpha。beta 是你通过投资整个市场，例如被动指数基金，所能获得的收益；由于管理者技能而额外获得的收益就是 alpha。 \\[0.65em]
Alternative beta & A kind of beta, but which requires active trading to achieve. So for example to earn the equity value premium, which is the return from being long low price:earnings (PE) and short high PE equities, you need to buy and sell the appropriate shares at the right time. Like beta, and unlike alpha, this kind of return does not require skill. \par 另类 beta：一种需要通过主动交易才能实现的 beta。例如，要赚取股票价值溢价，也就是做多低市盈率股票、做空高市盈率股票所得到的收益，你必须在适当时间买卖相应股票。和 beta 一样、与 alpha 不同，这类收益并不需要技能。 \\[0.65em]
Annualised cash volatility target & See volatility target. \par 年化现金波动率目标：见“波动率目标”。 \\[0.65em]
Asset allocating investor & An investor who usually does not forecast asset prices, but uses a systematic framework to create the best possible portfolio. See page ix and page 225. \par 资产配置型投资者：通常不预测资产价格，而是使用系统化框架来构建尽可能好的投资组合的投资者。见第 ix 页和第 225 页。 \\[0.65em]
Asset class & Name for a general kind of investable security, e.g. equities, bonds, commodity futures and commercial property are all asset classes. \par 资产类别：对某一大类可投资证券的统称，例如股票、债券、商品期货和商业地产都属于资产类别。 \\[0.65em]
Back-testing & A back-test is a historical simulation of what performance and behaviour would have been for a trading system, based on the data and prices that occurred in the past. See page 14. \par 回测：基于过去真实发生的数据和价格，对一个交易系统过去本会呈现的表现和行为进行历史模拟。见第 14 页。 \\[0.65em]
Beta & Used as a shorthand for getting returns from being exposed to the market generally, without having any particular skill; as distinct from alpha. Owning a passive index fund, such as an ETF linked to the French CAC 30 stock market, will give you beta returns for the relevant market. See also alternative beta. Also: to what extent a particular asset is exposed to the overall market. Low beta assets have less exposure to the overall market, and are safer, than high beta assets. \par Beta：常用来指无需特殊技能、仅通过暴露于整体市场而获得的收益，与 alpha 相对。持有一只被动指数基金，例如挂钩法国 CAC 30 股市的 ETF，就能获得相应市场的 beta 收益。另见“另类 beta”。beta 也可指某项资产相对于整体市场的暴露程度；低 beta 资产比高 beta 资产更少暴露于整体市场，也更安全。 \\[0.65em]
Block value & The increase in value of one instrument block when its price goes up by one percentage point. See page 154. \par 头寸块价值：当一个头寸块的价格上升 1 个百分点时，它的价值增加量。见第 154 页。 \\[0.65em]
Bootstrapping & A method of portfolio optimisation which takes into account the uncertainty in historical data, unlike single period optimisation. See page 75. \par 自助法优化：一种在投资组合优化时把历史数据中的不确定性纳入考虑的方法，与单期优化不同。见第 75 页。 \\[0.65em]
Calibration & Fine tuning a trading rule by looking at the performance or behaviour of variations during back-testing; a kind of fitting. See page 52. \par 校准：通过观察回测中不同变体的表现或行为，对交易规则进行细调的一种做法；属于拟合的一种。见第 52 页。 \\[0.65em]
Carry & A trading strategy where you profit from the difference between yields and funding costs if prices remain unchanged. A recommended carry trading rule is discussed in chapter seven, on page 119 and in appendix B, from page 285. See also FX Carry. \par Carry：一种在价格不变的情况下，通过收益率与融资成本之间差额来获利的交易策略。书中推荐的 carry 交易规则见第七章第 119 页，以及附录 B 第 285 页起。另见“外汇 carry”。 \\[0.65em]
Cash volatility target & See volatility target. \par 现金波动率目标：见“波动率目标”。 \\[0.65em]
Cognitive bias & A psychological flaw in human thought processes which results in poor decisions being made. See page 12. \par 认知偏差：人类思维过程中的一种心理缺陷，会导致糟糕决策。见第 12 页。 \\[0.65em]
Collective fund & A kind of fund which buys you a share in a portfolio of securities. These can be exchange traded funds, mutual funds in the US, or in the UK investment trusts and unit trusts. Collective funds can be active or passive. See page 106. \par 集合投资基金：一种让你持有一篮子证券投资组合份额的基金。它可以是 ETF、美国的共同基金，或英国的 investment trust 和 unit trust。集合投资基金既可以是主动型，也可以是被动型。见第 106 页。 \\[0.65em]
Combined forecast & The forecast you get for a particular instrument after combining the individual forecasts from multiple trading rules and multiplying by the forecast diversification multiplier. See page 125. \par 合成预测：对某一特定品种，把多条交易规则的单独预测组合起来，再乘以预测分散化乘数之后得到的预测值。见第 125 页。 \\[0.65em]
Correlation, Correlated & A measure of how two things co-move. Normally used in relation to daily returns from assets or profits from trading rules. A correlation of -1 indicates two things always move in opposite directions, +1 indicates they always move in the same direction and 0 means there is no linear relationship (uncorrelated). \par 相关性：衡量两个事物如何共同变动的指标。通常用于资产日收益或交易规则利润之间。相关性为 -1 表示两者总是反向运动，为 +1 表示总是同向运动，为 0 表示无线性关系，即不相关。 \\[0.65em]
Daily cash volatility target & See volatility target. \par 日现金波动率目标：见“波动率目标”。 \\[0.65em]
Data first & A kind of fitting where you specify a very general trading rule or model of price movement with many parameters and then use a statistical method to find the best choice of parameters from the available data. Opposite to ideas first. See page 26. \par 数据优先：一种拟合方法，先设定一个非常一般、包含许多参数的交易规则或价格运动模型，再用统计方法从现有数据中找出最优参数选择。与“想法优先”相对。见第 26 页。 \\[0.65em]
Data mining & The process of fitting to historical data in such a way that you will always find at least one apparently profitable trading rule. For example, if you run thousands of back-tests for different rules over a particular historical period, at least one will look great. Because they are highly tuned to fit the past, data mined trading rules usually perform badly in a future that is never the same. Also see over-fitting. \par 数据挖掘：以一种几乎总能在历史数据中找出至少一条“看起来赚钱”的交易规则的方式去拟合历史数据。例如，如果你对某一历史时期跑了几千次不同规则的回测，至少会有一条看起来非常漂亮。由于这类规则高度贴合过去，它们通常会在永远不再相同的未来表现很差。另见“过拟合”。 \\[0.65em]
Derivative & A way of benefiting from an asset going up or down in price without actually owning it. Futures, spread bets, options and contracts for difference are all examples of derivatives. Derivatives often provide more leverage than the underlying assets. \par 衍生品：一种无需真正持有资产本身，就能从其价格上涨或下跌中获利的工具。期货、点差投注、期权和差价合约都属于衍生品。衍生品通常能提供比标的资产更高的杠杆。 \\[0.65em]
Diversification multiplier & A portfolio will have a lower standard deviation of returns than its individual assets, assuming they are not perfectly correlated. This assumes the assets all have the same standard deviation of returns. The diversification multiplier multiplies portfolio returns to get them back to the same standard deviation as the underlying assets. It is always 1, for portfolios where all assets are perfectly correlated, or larger than 1. See instrument diversification multiplier and forecast diversification multiplier. See page 129. \par 分散化乘数：在资产并非完全相关、且各资产收益标准差相同的前提下，投资组合的收益标准差会低于单个资产。分散化乘数就是把组合收益重新放大回与底层资产相同标准差所需乘上的系数。若资产完全相关，该值恒为 1；否则大于 1。另见“品种分散化乘数”和“预测分散化乘数”。见第 129 页。 \\[0.65em]
Dynamic & A strategy where we actively buy and sell assets to express opinions on risk adjusted asset returns. As opposed to static. See page 38. \par 动态策略：一种通过主动买卖资产来表达对风险调整后资产收益看法的策略，与静态策略相对。见第 38 页。 \\[0.65em]
Equity value & A type of relative value trading in the equity market where you buy ‘cheap’ stocks with low price earnings (PE), high dividend yield or similar value characteristics, and sell ‘expensive’ stocks with high PE, low yield or similar. \par 股票价值策略：股票市场中的一种相对价值交易，做多低市盈率、高股息率或具有类似价值特征的“便宜”股票，同时做空高市盈率、低股息率或类似特征的“昂贵”股票。 \\[0.65em]
Equity volatility index & A measure of how volatile an equity market is expected to be, derived from option prices. Sometimes called the ‘fear index’. The key European index is the V2TX and in the US it’s the VIX. Futures can be traded on each of these indices. See skew concept box, page 34. \par 股票波动率指数：一种由期权价格推导出的、衡量股票市场预期波动程度的指标，有时也被称为“恐惧指数”。欧洲的关键指数是 V2TX，美国的是 VIX，这两个指数上都可以交易期货。见第 34 页关于 skew 的概念框。 \\[0.65em]
Exchange traded fund (ETF) & Usually a type of passive indexed, collective fund which can be traded and owned like a normal equity, and has low fees. Some ETFs implement dynamic strategies and are more expensive. \par 交易所交易基金（ETF）：通常是一种被动指数化的集合投资基金，可以像普通股票一样买卖和持有，费用较低。有些 ETF 也会实施动态策略，因此费用更高。 \\[0.65em]
Execution cost & Part of the cost of trading an instrument, equal to the difference between the mid-price and the traded price. Those who do not trade in large sizes can assume they will pay the bid (if selling) or offer (if buying), which implies it will cost them half the difference between the bid and the offer to trade. See page 179. \par 执行成本：交易某个品种时成本的一部分，等于中间价与实际成交价之间的差额。交易规模不大的交易者可以假设，卖出时成交在买价、买入时成交在卖价，这意味着执行成本大约等于买卖价差的一半。见第 179 页。 \\[0.65em]
Expanding window & A way of doing out of sample testing. You pretend you are at a point in the past and then use data from the beginning of your price history up to that point to choose trading rules or portfolio weights. You then test those rules or weights on the data for a year after that point. Then you move forward a year to the next point and repeat. See page 56. \par 扩张窗口：一种样本外测试方法。你假装自己处在过去某个时点，只使用从价格历史起点到该时点为止的数据来选择交易规则或组合权重；然后用该时点之后一年的数据进行测试，再向前推进一年重复这一过程。见第 56 页。 \\[0.65em]
Exponentially weighted moving average (EWMA) & A kind of moving average where you weight more recent observations more heavily. If \(P_t\) is the most recent data point then the EWMA is \((A \times P_t) + (A \times (1 - A) \times P_{t-1}) + (A \times (1 - A)^2 \times P_{t-2}) + \dots\). You can also calculate this recursively. If you have yesterday’s EWMA \(E_{t-1}\) then today’s EWMA is \((A \times P_t) + (E_{t-1} \times (1 - A))\). Used in the EWMAC trading rule and also in the estimation of price volatility. \par 指数加权移动平均（EWMA）：一种对较新观测赋予更高权重的移动平均。如果 \(P_t\) 是最新数据点，那么 EWMA 可写为 \((A \times P_t) + (A \times (1 - A) \times P_{t-1}) + (A \times (1 - A)^2 \times P_{t-2}) + \dots\)。也可以递推计算：若昨天的 EWMA 为 \(E_{t-1}\)，则今天的 EWMA 为 \((A \times P_t) + (E_{t-1} \times (1 - A))\)。用于 EWMAC 交易规则，也用于估计价格波动率。 \\[0.65em]
Exponentially weighted moving average crossover (EWMAC) & A trend following trading rule where you take a fast exponentially weighted moving average (EWMA) of the price with a short look-back, and a slow EWMA with a longer look-back. If the fast EWMA is above the slow then prices have been rising and you buy, and vice versa. See page 117 and page 282 in appendix B. \par 指数加权移动平均交叉（EWMAC）：一种趋势跟随交易规则，先对价格计算一条短回看期的快 EWMA 和一条长回看期的慢 EWMA。若快 EWMA 高于慢 EWMA，则说明价格近期在上涨，于是买入；反之亦然。见第 117 页及附录 B 第 282 页。 \\[0.65em]
Fitting & The process of picking the best trading rule. Usually this would involve running back-tests over historical data and finding the most profitable rule. See page 52. See also ideas first and data first, and over-fitted. \par 拟合：挑选最佳交易规则的过程。通常意味着在历史数据上运行回测，并找出最赚钱的规则。见第 52 页。另见“想法优先”、“数据优先”和“过拟合”。 \\[0.65em]
Forecast & An estimate of how much an instrument’s price will go up or down, translated into a numeric scale. The forecast can either be the discretionary forecast of a semi-automatic trader or the fixed forecast of the asset allocating investor; for staunch systems traders a single systematic trading rule or a combined forecast blending different trading rules. See page 110. \par 预测值：对某个品种价格将上涨或下跌多少的估计，并被转换成数值刻度。它可以是半自动交易者的主观预测、资产配置型投资者的固定预测；对坚定系统化交易者而言，也可以是一条系统化交易规则给出的预测，或多条规则混合后的合成预测。见第 110 页。 \\[0.65em]
Forecast diversification multiplier & The diversification multiplier required so that the combined forecast for a particular instrument has the expected average absolute value (recommended value: 10). See page 129. \par 预测分散化乘数：为了让某个品种的合成预测拥有期望的平均绝对值（推荐值为 10）所需要的分散化乘数。见第 129 页。 \\[0.65em]
Forecast scalar & A value you multiply a forecast by to ensure it has the right average absolute value (recommended value: 10). See page 297. \par 预测缩放因子：乘在预测值上的一个系数，用来确保预测值拥有正确的平均绝对值（推荐值为 10）。见第 297 页。 \\[0.65em]
Forecast weights & The weights used to combine forecasts from multiple trading rules into a single combined forecast for a particular instrument. See page 126. \par 预测权重：把多条交易规则的预测组合成某个品种单一合成预测时所使用的权重。见第 126 页。 \\[0.65em]
Foreign exchange (FX) carry & A foreign exchange (FX) trading strategy where you borrow in low interest currencies, exchange the money for a high interest currency, put it into the bank, and earn the difference. This strategy will fail if the high interest currency depreciates by more than the gap in interest rates. See also carry. \par 外汇 carry：一种外汇交易策略，即借入低利率货币，换成高利率货币后存入银行，赚取利差。如果高利率货币的贬值幅度超过利差，这个策略就会失败。另见“carry”。 \\[0.65em]
Framework & The process I describe in this book to translate forecasts into actual investment decisions. The framework includes components for combining trading rules and instruments, volatility targeting and position sizing. Together with any trading rules the framework makes up the trading system. See page 96. \par 框架：本书中我用来把预测值转化为实际投资决策的那整套流程。该框架包括交易规则与品种的组合、波动率目标设定和仓位规模设定等模块。框架与交易规则一起构成完整的交易系统。见第 96 页。 \\[0.65em]
Fundamental data & Non-price data used for forecasting, e.g. earnings of a firm, inflation rates or weather forecasts. Opposite of technical data. See page 43. \par 基本面数据：用于预测的非价格数据，例如公司盈利、通胀率或天气预报。与技术面数据相对。见第 43 页。 \\[0.65em]
Gaussian normal distribution & A statistical distribution which is shaped like a bell, often assumed for returns. If your returns are Gaussian normal then you will see returns one sigma or less around the average about 68\% of the time, and returns two sigma or less about 95\% of the time. See page 21. \par 高斯正态分布：一种钟形的统计分布，常被用作收益分布的假设。如果收益服从高斯正态分布，那么大约 68\% 的时候收益会落在均值上下 1 个 sigma 以内，约 95\% 的时候会落在 2 个 sigma 以内。见第 21 页。 \\[0.65em]
Half-Kelly & See Kelly criterion. \par 半 Kelly：见“Kelly 准则”。 \\[0.65em]
Handcrafted optimisation & A method of portfolio optimisation where you set weights by hand, grouping assets together and using only estimates of correlations (and optionally, Sharpe ratios). See page 78. \par 手工构造优化：一种投资组合优化方法，通过把资产分组并用相关性估计（以及可选的夏普比率）来手工设定权重。见第 78 页。 \\[0.65em]
Ideas first & A type of fitting where you conceive an idea, choose a specific trading rule and then back-test the rule on historical data. Opposite to data first. See page 26. \par 想法优先：一种拟合方式，先形成一个想法，再选择具体交易规则，然后在历史数据上回测该规则。与“数据优先”相对。见第 26 页。 \\[0.65em]
In sample & Using historical data to select trading rules or portfolio weights and then testing those rules on the same data, which means including future information. A bad thing. Opposite to out of sample. See page 54. \par 样本内：用历史数据来选择交易规则或组合权重，然后又在同一批数据上测试这些规则，这意味着包含了未来信息。不是好事。与“样本外”相对。见第 54 页。 \\[0.65em]
Index tracker & See passive fund. \par 指数跟踪基金：见“被动基金”。 \\[0.65em]
Instrument & Something you trade or invest in, e.g. an equity like Apple stock, an oil futures contract or a spread bet on the EUR/USD FX rate. See page 101. \par 品种：你用来交易或投资的对象，例如苹果股票这样的权益资产、原油期货合约，或 EUR/USD 汇率上的点差投注。见第 101 页。 \\[0.65em]
Instrument block & ‘One’ of an instrument – the minimum discrete unit you can economically trade an instrument in, which might be 100 shares, one futures contract or £1 a point on a spread bet. See page 154. \par 头寸块：某个品种在经济上可交易的最小离散单位，比如 100 股、一张期货合约，或点差投注中的每点 1 英镑。见第 154 页。 \\[0.65em]
Instrument currency volatility & The expected daily standard deviation of returns of owning one instrument block, measured in the currency of the instrument. Equal to price volatility multiplied by block value. See page 158. \par 品种货币波动率：持有一个头寸块时，以该品种计价货币衡量的预期日收益标准差。等于价格波动率乘以头寸块价值。见第 158 页。 \\[0.65em]
Instrument diversification multiplier & A type of diversification multiplier that accounts for the diversification across the returns of trading subsystems. See page 169. \par 品种分散化乘数：一种分散化乘数，用来反映不同交易子系统收益之间的分散化效果。见第 169 页。 \\[0.65em]
Instrument value volatility & The expected daily standard deviation of returns of owning one instrument block, measured in the same currency as the trader or investor’s trading capital. Equal to instrument currency volatility multiplied by the relevant exchange rate. See page 158. \par 品种价值波动率：持有一个头寸块时，以交易者或投资者交易资本所使用货币衡量的预期日收益标准差。等于品种货币波动率乘以相应汇率。见第 158 页。 \\[0.65em]
Instrument weights & In my framework the weights that different trading subsystems, each trading a single instrument, have in a portfolio of subsystems. See page 165. \par 品种权重：在我的框架中，不同交易子系统在子系统组合中的权重；每个交易子系统只交易一个品种。见第 165 页。 \\[0.65em]
Kelly criterion & A way of determining the percentage volatility target you should use, given your expectation of Sharpe ratio (SR). A SR of 0.5 implies you should use a risk percentage of 50\%. However it is safer to use Half-Kelly and set the risk to half what Kelly requires, i.e. for an SR of 0.5 set risk to 25\%. See page 143. \par Kelly 准则：一种根据你预期的夏普比率（SR）来确定应使用多少百分比波动率目标的方法。若 SR 为 0.5，则意味着风险比例应为 50\%。不过使用 Half-Kelly 更安全，也就是把 Kelly 要求的风险减半；例如 SR 为 0.5 时，风险设为 25\%。见第 143 页。 \\[0.65em]
Law of active management & An equation which states that the information ratio of a trading strategy is proportional to the square root of the number of independent bets made each year. See page 42. \par 主动管理法则：一条公式，指出交易策略的信息比率与每年进行的独立下注次数的平方根成正比。见第 42 页。 \\[0.65em]
Leverage & Borrowing to invest, either in an explicit way or by using a derivative such as a future or spread bet where your exposure is greater than your initial cash payment. \par 杠杆：借钱进行投资，可以是显性的借款，也可以是通过期货或点差投注等衍生品实现，使你的风险暴露大于最初现金支付。 \\[0.65em]
Liquidity & How easy it is to trade quickly in size without changing the price significantly. Shares in blue chip companies are liquid because you can usually sell \$1 million worth within minutes at close to the prevailing price. A \$1 million house cannot be sold in a single day without giving buyers a substantial discount from what it could achieve given more time. The blue chip shares are liquid; the house is not. See page 35. \par 流动性：在不显著影响价格的前提下，快速、大规模交易某资产的容易程度。蓝筹股具有流动性，因为你通常可以在几分钟内以接近当前价格卖出价值 100 万美元的股票；而一套价值 100 万美元的房子若想一天内卖掉，就必须大幅打折。蓝筹股有流动性，房子则没有。见第 35 页。 \\[0.65em]
Mean reversion & Any trading strategy where you assume asset prices will revert to an equilibrium or fair value. For example, you might think 1.70 is a fair value for the price of the GBP/USD FX rate. If the rate goes below 1.70 you’ll buy pounds and sell dollars, and if it goes higher you’d sell pounds (and buy dollars). See also relative value. \par 均值回归：任何假设资产价格会回归某个均衡值或公允价值的交易策略。例如，你可能认为 GBP/USD 的公允价格是 1.70；若汇率跌到 1.70 以下，你就买英镑卖美元；若高于 1.70，你就卖英镑买美元。另见“相对价值”。 \\[0.65em]
Merger arbitrage & A trading strategy, where if a merger or takeover is occurring you buy the company to be acquired, and usually short sell the acquirer as a hedge. Normally this is done at a discount to the takeover price, the gap reflecting uncertainty about whether a deal will go ahead. Profits can be made by selecting deals where the gap is too large given the level of uncertainty. This is a negative skew style of trading, because occasional large losses are made when deals fall through. \par 并购套利：一种在并购或收购发生时买入被收购公司、通常同时做空收购方以作对冲的交易策略。通常买入价格会低于收购价，这个价差反映了交易能否完成的不确定性。若你能挑出在给定不确定性下价差过大的交易，就能获利。这是一种负偏度交易风格，因为一旦交易告吹，就会偶尔出现较大亏损。 \\[0.65em]
Momentum & When asset prices go up after previous rises; or go down after falls. Trend following rules try and profit from momentum. \par 动量：资产价格在之前上涨后继续上涨，或在下跌后继续下跌的现象。趋势跟随规则试图从动量中获利。 \\[0.65em]
Moving average & A weighted average of a value over the last N points in time, where N is the look-back window. The weighting can either be a simple average, with all points equally weighted, or an exponentially weighted moving average, with more recent points given a higher weight. Used in the estimation of price volatility. \par 移动平均：对过去 N 个时间点数值做加权平均，其中 N 是回看窗口。权重既可以是简单平均，即所有点权重相同，也可以是指数加权移动平均，即较近的数据点权重更高。用于估计价格波动率。 \\[0.65em]
Open interest & The number of futures contracts outstanding for a particular instrument. A measure of how actively an instrument is traded. \par 未平仓量：某个特定品种尚未平仓的期货合约数量，是衡量该品种交易活跃程度的指标。 \\[0.65em]
Out of sample & Using historical data to select trading rules or portfolio weights without looking into the future. A good thing, in contrast to in sample. See page 54. \par 样本外：在不看未来信息的前提下，用历史数据选择交易规则或组合权重。与样本内相对，是好事。见第 54 页。 \\[0.65em]
Over-fitted & When something is fitted in such a way that it matches the past data too well and is unlikely to perform well when actually traded. Sometimes referred to as curve fitting. Also see data mining. Refer to page 52. \par 过拟合：某个模型或规则被拟合得过于贴合过去数据，以至于在真实交易时不太可能表现良好。有时也称为曲线拟合。另见“数据挖掘”。参见第 52 页。 \\[0.65em]
Passive fund, Passive index, Passive management & A type of collective fund which is invested and managed passively in a portfolio of securities according to some predefined weights. Usually these funds are index trackers which follow indices such as the FTSE 100 or S\&P 500. Passive funds are a cheap way to get broad market, or beta, exposure. As distinct from active funds. See page 106. \par 被动基金、被动指数、被动管理：一种集合投资基金，按照预先设定好的权重被动地投资和管理一篮子证券。它们通常是跟踪 FTSE 100 或 S\&P 500 这类指数的指数基金。被动基金是获得广义市场，也就是 beta，暴露的一种低成本方式。与主动基金相对。见第 106 页。 \\[0.65em]
Percentage volatility target & The target expected annualised percentage risk your trading system is exposed to. When multiplied by trading capital we get the annualised cash volatility target. See also volatility target and page 138. \par 百分比波动率目标：你的交易系统所暴露的目标年化百分比风险。将其乘以交易资本，就得到年化现金波动率目标。另见“波动率目标”和第 138 页。 \\[0.65em]
Portfolio optimisation & The process of finding the optimal (best) portfolio by deciding in which portfolio weights to hold your assets. Normally a portfolio consists of positions in various assets, e.g. stocks and bonds. In my framework portfolios can be of (a) multiple trading rules for which we find forecast weights, and for (b) trading subsystems, one per instrument, for which we optimise instrument weights. You should use handcrafting or bootstrapping to calculate portfolio weights; I do not recommend single period optimisation. See page 70. \par 投资组合优化：通过决定各资产应持有的组合权重，来寻找最优投资组合的过程。通常投资组合由股票、债券等多种资产头寸组成。在我的框架里，组合既可以是多条交易规则组成的组合，此时我们寻找的是预测权重；也可以是按品种划分的多个交易子系统组成的组合，此时我们优化的是品种权重。我建议用手工构造法或自助法来计算组合权重，不推荐单期优化。见第 70 页。 \\[0.65em]
Portfolio weighted position & The position you hold in an instrument at the trading system level. Because a trading system consists of a portfolio of trading subsystems, one for each instrument, this is equal to the instrument’s subsystem position multiplied by the instrument weight and the instrument diversification multiplier. See page 173. \par 组合加权仓位：在交易系统层面上，你在某个品种上持有的仓位。由于交易系统是由多个交易子系统组成的组合，每个品种对应一个子系统，所以该仓位等于该品种子系统仓位乘以品种权重和品种分散化乘数。见第 173 页。 \\[0.65em]
Portfolio weights & In general the weights that assets have in a portfolio. See instrument weights and forecast weights. \par 组合权重：一般意义上资产在投资组合中的权重。另见“品种权重”和“预测权重”。 \\[0.65em]
Position inertia & If the current position is less than 10\% away from the rounded target position then to reduce costs you shouldn’t trade. See page 174. \par 仓位惯性：如果当前仓位与四舍五入后的目标仓位相差不到 10\%，那么为了降低成本，你就不应交易。见第 174 页。 \\[0.65em]
Predictable risk & The component of volatility in asset prices or portfolio returns which you can predict using estimates of standard deviation and correlation. See unpredictable risk and page 39. \par 可预测风险：资产价格或组合收益波动中，能够通过标准差和相关性估计来预测的那一部分。另见“不可预测风险”和第 39 页。 \\[0.65em]
Price volatility & The expected standard deviation of instrument price daily returns, in percentage points per day. See page 155. \par 价格波动率：某个品种价格日收益的预期标准差，以每日百分比点计。见第 155 页。 \\[0.65em]
Relative value & Any trading strategy where you buy something that is cheap, and sell something similar that is expensive as a hedge. For example, see equity value. Often a negative skew strategy, as we suffer occasional large losses. \par 相对价值：任何做多便宜资产、同时做空相似但昂贵资产以作对冲的交易策略。例如“股票价值策略”。它通常是一种负偏度策略，因为会偶尔遭受较大亏损。 \\[0.65em]
Risk parity & A type of passive fund where assets are held in equal shares according to expected risk. See page 38. \par 风险平价：一种被动基金形式，按照预期风险把资产以相等风险份额持有。见第 38 页。 \\[0.65em]
Risk premium & The reward you get from investing in something that carries some risk. For example, equities are riskier than bonds so we should expect higher returns to compensate us. See page 31. \par 风险溢价：投资于带有某种风险的资产时所获得的补偿。例如，股票比债券风险更高，因此我们应当预期股票提供更高收益作为补偿。见第 31 页。 \\[0.65em]
Rolling window & A method of out of sample testing. You pretend you are at a point in the past and then use data from the last N years (where N can be any value) up to that point to choose trading rules or portfolio weights. You then test those rules or weights on data for a year after that point. Then you move forward a year to the next point and repeat. See page 56. \par 滚动窗口：一种样本外测试方法。你假装自己身处过去某个时点，只使用此前最近 N 年的数据来选择交易规则或组合权重；随后用该时点之后一年的数据测试这些规则或权重，再向前推进一年重复这一过程。见第 56 页。 \\[0.65em]
Semi-Automatic Trader & A trader who makes their own forecasts about price movements in a discretionary fashion, but then incorporates them into a systematic framework. See page ix and page 209. \par 半自动交易者：以主观方式自行判断价格走势，但随后把这些判断纳入系统化框架中执行的交易者。见第 ix 页和第 209 页。 \\[0.65em]
Sharpe ratio (SR) & A measure of how profitable a trading strategy is, with returns adjusted for risk. Formally it is the mean return over some time period divided by the standard deviation of returns over the same time period. In this book I normally use the annualised Sharpe ratio – annualised returns divided by annualised standard deviation. If we’re not using derivatives the ‘risk free’ interest rate should be deducted from annualised returns before calculating the Sharpe ratio. See page 32. \par 夏普比率（SR）：衡量交易策略在进行风险调整后有多盈利的指标。形式上，它等于某一时期平均收益除以同一时期收益标准差。本书里我通常使用年化夏普比率，即年化收益除以年化标准差。如果不用衍生品，那么在计算夏普比率之前，应先从年化收益中扣除无风险利率。见第 32 页。 \\[0.65em]
Short selling & A way of benefiting from an asset’s fall in price, usually used for equity trading. You would usually borrow the stock with a loan agreement before selling it. On closing the trade the stock is returned and you’ll profit if the price has fallen since inception. \par 卖空：一种从资产价格下跌中获利的方式，通常用于股票交易。你一般会先通过借券协议借来股票再卖出；平仓时再把股票归还，如果价格较开仓时下跌，你就能获利。 \\[0.65em]
Sigma & Shorthand for one unit of standard deviation. See page 21. \par Sigma：一个标准差单位的简写。见第 21 页。 \\[0.65em]
Single period optimisation & The classic way of performing a portfolio optimisation: using a single average over historical data of asset return means, standard deviations and correlations; these estimates are then used to do a single optimisation. \par 单期优化：进行投资组合优化的经典方法，即用历史数据的单一平均值来估计资产收益均值、标准差和相关性，再用这些估计值做一次优化。 \\[0.65em]
Skew & A measure of how symmetric the returns from an asset or trading rule are. Positive skew means you get more losses and fewer gains. But the average loss is smaller than the average gain. Negative skew means you get more gains and fewer losses, and average losses are larger than average gains. See page 32. \par 偏度：衡量资产或交易规则收益分布有多对称的指标。正偏度意味着亏损次数更多、盈利次数更少，但平均亏损小于平均盈利；负偏度则意味着盈利次数更多、亏损次数更少，但平均亏损大于平均盈利。见第 32 页。 \\[0.65em]
Speed limit & A limit on how quickly you should trade an instrument. The speed limit will depend on the standardised cost of the instrument, and how much you are prepared to pay in Sharpe ratio (SR) units per year in costs. I recommend paying no more than 0.13 SR units if you are a staunch systems trader and 0.08 SR units for semi-automatic traders and asset-allocating investors. See page 187. \par 速度上限：你在某个品种上应当交易多快的上限。它取决于该品种的标准化成本，以及你愿意每年用多少夏普比率单位来支付成本。我的建议是：坚定系统化交易者每年支付的成本不超过 0.13 SR 单位，半自动交易者和资产配置型投资者不超过 0.08 SR 单位。见第 187 页。 \\[0.65em]
Standard deviation & A measure of how dispersed some data is around its average value. If your data points are \(x_1, x_2, \ldots, x_n\) then the average \(x^* = (x_1 + x_2 + \ldots + x_n) \div n\). The standard deviation is \(\sqrt{(1 \div n)[(x_1 - x^*)^2 + (x_2 - x^*)^2 + \ldots + (x_n - x^*)^2]}\). Often applied to daily returns in prices or trading system profits. An approximate way to calculate an annualised standard deviation of returns is to multiply the daily standard deviation by 16, the square root of 256, approximately the number of trading days in one year. See page 21. \par 标准差：衡量一组数据围绕其平均值分散程度的指标。若数据点为 \(x_1, x_2, \ldots, x_n\)，则平均值 \(x^* = (x_1 + x_2 + \ldots + x_n) \div n\)。标准差为 \(\sqrt{(1 \div n)[(x_1 - x^*)^2 + (x_2 - x^*)^2 + \ldots + (x_n - x^*)^2]}\)。它常被应用在价格日收益或交易系统利润上。要近似年化日收益标准差，可以乘以 16，也就是 256 的平方根，约等于一年交易日数的平方根。见第 21 页。 \\[0.65em]
Standardised cost & A volatility standardised method of measuring costs which is comparable across instruments. To calculate, add up all the costs of trading \(C\), including execution costs, fees and taxes. Then double to get the cost of a round trip, a buy and a sell, and divide by the annualised instrument cash volatility \(ICV\) of the instrument. The formula is \(2 \times C \div (16 \times ICV)\). See page 181. \par 标准化成本：一种经过波动率标准化、可跨品种比较的成本度量方式。计算时先把所有交易成本 \(C\) 加总，包括执行成本、费用和税费；然后乘以 2 得到一次往返交易，也就是一次买入和一次卖出的总成本，再除以该品种的年化品种现金波动率 \(ICV\)。公式为 \(2 \times C \div (16 \times ICV)\)。见第 181 页。 \\[0.65em]
Static & A strategy where you do not actively trade the assets you are investing in because of changes in expected risk adjusted returns. As opposed to dynamic. See page 38. \par 静态策略：一种不会因为预期风险调整后收益发生变化而主动交易所投资资产的策略，与动态策略相对。见第 38 页。 \\[0.65em]
Staunch Systems Trader & A trader who uses one or more systematic trading rules and incorporates them into a systematic framework. See page x and page 245. \par 坚定系统化交易者：使用一条或多条系统化交易规则，并把它们纳入系统化框架中的交易者。见第 x 页和第 245 页。 \\[0.65em]
Subsystem position & The position nominally held by a single instrument’s trading subsystem, trading only one instrument with the entire trading capital, given a forecast of returns. Equal to the volatility scalar multiplied by forecast, divided by 10. See page 159. \par 子系统仓位：某个单一品种交易子系统在“用全部交易资本只交易这一个品种”的假想条件下、基于给定收益预测而持有的名义仓位。等于波动率缩放因子乘以预测值再除以 10。见第 159 页。 \\[0.65em]
Survivorship bias & A problem where assets that disappeared are missing from historic data sets, such as shares in firms which went bankrupt. Because we don’t see the losses from these instruments we are likely to overestimate how profitable investing really was in the past. \par 幸存者偏差：一种问题，即那些后来消失的资产，比如破产公司的股票，缺失在历史数据集中。由于这些资产带来的亏损没有被看见，我们就容易高估过去投资实际上有多赚钱。 \\[0.65em]
Technical & A kind of trading rule that uses only price data to predict prices, such as trend following; no fundamental data is used. See page 43. \par 技术面：一种只使用价格数据来预测价格的交易规则，例如趋势跟随；不使用任何基本面数据。见第 43 页。 \\[0.65em]
Trading capital & The amount of capital at risk in your trading system. Also see volatility target. See page 138. \par 交易资本：在你的交易系统中承担风险的资本数量。另见“波动率目标”。见第 138 页。 \\[0.65em]
Trading rules & A systematic rule used to predict whether the price of an instrument will go up or down, and by how much. When one or more trading rules are put into a framework they form a complete trading system for trading systematically. See page 109. \par 交易规则：一种系统化规则，用来预测某个品种价格会上涨还是下跌，以及幅度有多大。当一条或多条交易规则被放入框架中时，它们就构成了完整的系统化交易系统。见第 109 页。 \\[0.65em]
Trading subsystem & A trading subsystem is a notional part of a larger trading system which trades a single instrument. See page 103. \par 交易子系统：更大交易系统中的一个名义组成部分，只负责交易单一品种。见第 103 页。 \\[0.65em]
Trading system & A set of one or more trading rules within a framework creates a complete trading system; it makes forecasts about instrument price movements and translates these into positions and hence trades suitable for a particular level of trading capital. For asset allocating investors and staunch systems traders the trading system is made up of a portfolio of trading subsystems, one per instrument, for a fixed set of instruments and trading rules (although asset allocating investors use only one rule with a constant forecast). For semi-automatic traders the system is made up of discretionary forecasts for an ad hoc group of instruments, so the make-up of the group of trading subsystems will vary over time. See page 98. \par 交易系统：在框架中放入一条或多条交易规则，就形成了完整的交易系统；它对品种价格运动做出预测，并把预测转化为适合特定交易资本水平的仓位和交易。对资产配置型投资者和坚定系统化交易者而言，交易系统由多个交易子系统组成的组合构成，每个品种对应一个子系统，并针对固定的一组品种和交易规则运作（不过资产配置型投资者只使用一条恒定预测的规则）。对半自动交易者而言，系统则由一组临时性的主观预测组成，因此交易子系统集合会随时间变化。见第 98 页。 \\[0.65em]
Trend following & A technical trading strategy which tries to capture momentum in price; you buy things going up and sell things that have fallen in price. Tends to have positive skew. \par 趋势跟随：一种试图捕捉价格动量的技术交易策略；买入正在上涨的东西，卖出已经下跌的东西。它通常具有正偏度。 \\[0.65em]
Turnover & A way of measuring trading speed. Turnover is measured in round trips per year, where a round trip consists of a buy and a sell of an average sized position. So a trading system with a turnover of 10 units is expected to do ten buys and ten sells in a year. See page 184. \par 换手率：衡量交易速度的一种方式。它以每年完成多少次往返交易来度量，而一次往返交易由对平均仓位做一次买入和一次卖出组成。因此，换手率为 10 的交易系统，预期一年会做 10 次买入和 10 次卖出。见第 184 页。 \\[0.65em]
Unpredictable risk & The component of volatility in asset prices or portfolio returns which you didn’t or couldn’t predict in advance. This could be because standard deviations or correlations changed, or because your risk model was wrong. See predictable risk and page 39. \par 不可预测风险：资产价格或组合收益波动中，那部分你事先没有预测到、或者根本无法预测到的成分。这可能是因为标准差或相关性发生了变化，也可能是因为你的风险模型本身就错了。另见“可预测风险”和第 39 页。 \\[0.65em]
Variation & A variation on a trading rule has different parameter value(s) but is otherwise the same. For example ‘buy in the range 1.5 to 2.0’ is a variation of the general rule of ‘buy in the range X to Y’. See page 53. \par 变体：交易规则的某个变体，除了参数值不同之外，其余部分都相同。例如“在 1.5 到 2.0 区间买入”就是一般规则“在 X 到 Y 区间买入”的一个变体。见第 53 页。 \\[0.65em]
Volatility & A shorthand term for standard deviation. \par 波动率：标准差的简写说法。 \\[0.65em]
Volatility look-back period & Period of recent time used to calculate a standard deviation of returns. Used to calculate the price volatility and in trading rules like EWMAC. \par 波动率回看期：用于计算收益标准差的一段最近时间窗口。它既用于计算价格波动率，也用于 EWMAC 之类的交易规则。 \\[0.65em]
Volatility scalar & A parameter which accounts for the difference in the risk you want for your portfolio (volatility target) relative to the risk of an instrument. It’s the number of instrument blocks you would hold if you invested your entire trading capital into one instrument, with a forecast of +10. Equal to daily cash volatility target divided by instrument value volatility. See page 159. \par 波动率缩放因子：一个用来处理“你希望组合承担的风险（波动率目标）”与“单个品种本身风险”之间差异的参数。它表示：如果你把全部交易资本都投入到某一个品种，且该品种预测值为 +10，那么你应该持有多少个头寸块。它等于日现金波动率目标除以品种价值波动率。见第 159 页。 \\[0.65em]
Volatility standardisation & Adjusting returns or costs so they have the same expected risk. See page 40. \par 波动率标准化：调整收益或成本，使它们拥有相同的预期风险。见第 40 页。 \\[0.65em]
Volatility target & The target expected standard deviation of returns for the entire trading system. Can be expressed in cash terms, in the same currency as trading capital, either as the annualised cash volatility target or daily cash volatility target. The daily cash volatility target is the annualised cash volatility target divided by ‘the square root of time’, which assuming a 256 business day year is 16. It can also be expressed as a percentage volatility target, as a percentage of trading capital. The cash targets are equal to trading capital multiplied by the percentage volatility target. See page 137. \par 波动率目标：整个交易系统目标上的预期收益标准差。它可以用现金金额表示，使用与交易资本相同的货币，具体可以写成年化现金波动率目标或日现金波动率目标。日现金波动率目标等于年化现金波动率目标除以“时间平方根”；若一年按 256 个交易日计，这个值就是 16。它也可以写成百分比波动率目标，即占交易资本的百分比。现金目标等于交易资本乘以百分比波动率目标。见第 137 页。 \\[0.65em]
\end{longtable}
}
\end{chinesetableblock}
